\chapter{Orbitales frontières et réactivité}\label{ch:b2:frontier-orbitals}

Un flacon de cyclopentadiène oublié sur une étagère se transforme lentement en son dimère~:
deux molécules du même \omterm{def:b2:frontier-orbitals:diels-alder}{diène} s'unissent en un composé bicyclique ponté, et
un seul des stéréoisomères possibles se forme. Rien, dans les flèches courbes du
volume de première année, ne permet de prévoir lequel. Deux orbitales le permettent~: la plus haute orbitale
occupée d'une molécule et la plus basse orbitale vacante de l'autre. Ce
chapitre fait de l'interaction à deux niveaux du \cref{ch:b2:diatomic-mos} un
outil pour la réactivité --- quel atome un nucléophile attaque, de quel côté, et
pourquoi un \omterm{def:b2:frontier-orbitals:diels-alder}{diène} et un alcène s'unissent en un cycle en une seule étape.

\begin{recall}
\Cref{ch:b2:diatomic-mos}~: deux orbitales en interaction, répulsion à quatre électrons.
\Cref{ch:b2:huckel}~: niveaux et coefficients de Hückel des systèmes conjugués.
Le volume de première année~: nucléophiles et électrophiles, flèches courbes, contrôle
cinétique, réactions stéréospécifiques, stéréosélectives et régiosélectives,
additions électrophiles sur les alcènes, substitution nucléophile.
\end{recall}

\section{Deux orbitales, de nouveau}

Lorsque deux molécules s'approchent, leurs orbitales se recouvrent et interagissent. L'essentiel de
l'interaction est faible, les orbitales étant éloignées en énergie~; une forme perturbative du
résultat à deux niveaux suffit alors.

\begin{theorem}[Interaction faible de deux niveaux]\label{thm:b2:frontier-orbitals:perturbation}
Deux orbitales d'énergies $\varepsilon_1 < \varepsilon_2$, couplées par une \omterm{def:b2:diatomic-mos:integrals}{intégrale de résonance}
$\beta$ telle que $|\beta| \ll \Delta\varepsilon = \varepsilon_2 - \varepsilon_1$ (recouvrement négligé), sont
déplacées en
\[
  \varepsilon_1 - \frac{\beta^2}{\Delta\varepsilon}, \qquad \varepsilon_2 + \frac{\beta^2}{\Delta\varepsilon}
\]
au premier ordre en $\beta^2/\Delta\varepsilon^2$~: le niveau inférieur descend et le niveau supérieur monte de la même
quantité, d'autant plus petite que l'écart est grand.
\end{theorem}

\begin{proof}
D'après le \cref{thm:b2:diatomic-mos:heteronuclear}, les niveaux exacts sont $\bar\varepsilon \mp
\sqrt{\Delta\varepsilon^2/4 + \beta^2}$. Avec $u = 4\beta^2/\Delta\varepsilon^2 \ll 1$, $\sqrt{\Delta\varepsilon^2/4
+ \beta^2} = \frac{\Delta\varepsilon}{2}\sqrt{1 + u} \approx \frac{\Delta\varepsilon}{2}(1 + u/2) =
\frac{\Delta\varepsilon}{2} + \frac{\beta^2}{\Delta\varepsilon}$, et $\bar\varepsilon \mp \Delta\varepsilon/2 = \varepsilon_1,
\varepsilon_2$.
\end{proof}

\begin{proposition}[Deux électrons attirent, quatre repoussent]\label{prop:b2:frontier-orbitals:four-electron}
Si les deux orbitales en interaction contiennent deux électrons en tout, le système est stabilisé
d'environ $2\beta^2/\Delta\varepsilon$. Si elles en contiennent quatre, le système est déstabilisé.
\end{proposition}

\begin{proof}
Deux électrons vont dans le niveau abaissé~: gain $2\beta^2/\Delta\varepsilon$. Avec quatre, le niveau supérieur
est rempli lui aussi~; à l'ordre du \cref{thm:b2:frontier-orbitals:perturbation}, les deux
déplacements se compensent, et en tenant compte du recouvrement le niveau supérieur monte plus que l'inférieur ne
descend (\cref{prop:b2:diatomic-mos:asymmetry})~: une répulsion nette.
\end{proof}

\begin{omfigure}
\begin{tikzpicture}[font=\small]
  \begin{scope}
    \pic (a) at (0,0) {omlevel={ud}{0.7}};
    \pic (b) at (3,1.4) {omlevel={empty}{0.7}};
    \pic (lo) at (1.5,-0.5) {omlevel={ud}{0.7}};
    \pic (hi) at (1.5,1.9) {omlevel={empty}{0.7}};
    \draw[omcorr, densely dashed] (a-r) -- (lo-l) (a-r) -- (hi-l) (b-l) -- (lo-r) (b-l) -- (hi-r);
    \node[below=6pt] at (1.5,-0.6) {deux électrons~: stabilisation};
    \node[left] at (a-l) {pleine};
    \node[right] at (b-r) {vacante};
  \end{scope}
  \begin{scope}[xshift=7cm]
    \pic (a) at (0,0) {omlevel={ud}{0.7}};
    \pic (b) at (3,0.6) {omlevel={ud}{0.7}};
    \pic (lo) at (1.5,-0.5) {omlevel={ud}{0.7}};
    \pic (hi) at (1.5,1.25) {omlevel={ud}{0.7}};
    \draw[omcorr, densely dashed] (a-r) -- (lo-l) (a-r) -- (hi-l) (b-l) -- (lo-r) (b-l) -- (hi-r);
    \node[below=6pt] at (1.5,-0.6) {quatre électrons~: répulsion};
    \node[left] at (a-l) {pleine};
    \node[right] at (b-r) {pleine};
  \end{scope}
\end{tikzpicture}

{\small Interaction de deux orbitales de molécules voisines. Une orbitale pleine
face à une orbitale vacante donne une stabilisation nette~; deux orbitales pleines se repoussent, le
niveau supérieur montant plus que l'inférieur ne descend. Cette répulsion est l'origine
de la gêne stérique entre molécules.}
\end{omfigure}

\section{Les orbitales frontières}

\begin{definition}[Orbitales frontières]\label{def:b2:frontier-orbitals:frontier}
La \emph{HOMO}\index{HOMO} (de l'anglais \textit{highest occupied molecular orbital}~: la plus haute \omterm{def:b2:diatomic-mos:lcao}{orbitale
moléculaire} occupée, notée aussi HO) et la \emph{LUMO}\index{LUMO} (\textit{lowest unoccupied molecular orbital}~: la plus basse
\omterm{def:b2:diatomic-mos:lcao}{orbitale moléculaire} vacante, notée aussi BV) d'une molécule sont ses \emph{orbitales frontières}\index{orbitales frontieres@orbitales frontières}.
\end{definition}

\begin{proposition}[L'interaction dominante]\label{prop:b2:frontier-orbitals:dominant}
Entre deux molécules à couches complètes, les interactions stabilisantes qui comptent le plus
sont celles entre la \omterm{def:b2:frontier-orbitals:frontier}{HOMO} de l'une et la \omterm{def:b2:frontier-orbitals:frontier}{LUMO} de l'autre~; des deux paires de ce type, celle
dont l'écart est le plus petit l'emporte. Un nucléophile réagit par sa \omterm{def:b2:frontier-orbitals:frontier}{HOMO}, un
électrophile par sa \omterm{def:b2:frontier-orbitals:frontier}{LUMO}.
\end{proposition}

\begin{proof}[Argument]
Les paires pleine--pleine se repoussent (\cref{prop:b2:frontier-orbitals:four-electron}) et les paires
vacante--vacante ne contiennent aucun électron~: seules les paires pleine--vacante stabilisent, chacune de $2\beta^2/\Delta\varepsilon$.
La \omterm{def:b2:frontier-orbitals:frontier}{HOMO} est le plus haut niveau rempli et la \omterm{def:b2:frontier-orbitals:frontier}{LUMO} le plus bas niveau vacant~: une
paire HOMO--LUMO a donc le plus petit écart parmi elles~; d'après le théorème, le plus petit écart
donne le terme le plus grand.
\end{proof}

\begin{definition}[Contrôle de charge et contrôle orbitalaire]\label{def:b2:frontier-orbitals:control}
Une réaction est sous \emph{contrôle de charge}\index{controle de charge@contrôle de charge} lorsque l'attraction
entre les charges des deux partenaires décide de l'endroit où ils réagissent, et sous
\emph{contrôle orbitalaire}\index{controle orbitalaire@contrôle orbitalaire} lorsque c'est le recouvrement des \omterm{def:b2:frontier-orbitals:frontier}{orbitales frontières}
qui en décide, l'attaque ayant lieu là où leurs coefficients sont les plus grands.
\end{definition}

\begin{definition}[Espèces dures et molles]\label{def:b2:frontier-orbitals:hard-soft}
Un \emph{nucléophile dur}\index{nucleophile dur@nucléophile dur} ou un \emph{électrophile dur}\index{electrophile dur@électrophile dur}
est petit, porte une charge concentrée et possède
une \omterm{def:b2:frontier-orbitals:frontier}{HOMO} (respectivement une \omterm{def:b2:frontier-orbitals:frontier}{LUMO}) éloignée de celle de ses partenaires~: il réagit sous \omterm{def:b2:frontier-orbitals:control}{contrôle de
charge}. Un \emph{nucléophile mou}\index{nucleophile mou@nucléophile mou} ou un \emph{électrophile mou}\index{electrophile mou@électrophile mou}
est gros et polarisable, avec une \omterm{def:b2:frontier-orbitals:frontier}{HOMO} haute
(respectivement une \omterm{def:b2:frontier-orbitals:frontier}{LUMO} basse)~: il réagit sous \omterm{def:b2:frontier-orbitals:control}{contrôle orbitalaire}. Le dur réagit de préférence
avec le dur, le mou avec le mou.
\end{definition}

L'ion hydroxyde et l'ion fluorure sont des \omterm{def:b2:frontier-orbitals:hard-soft}{nucléophiles durs}, l'iodure et les thiolates
des \omterm{def:b2:frontier-orbitals:hard-soft}{nucléophiles mous}~; un proton et un carbocation sont des \omterm{def:b2:frontier-orbitals:hard-soft}{électrophiles durs}, le carbone d'un
halogénoalcane un \omterm{def:b2:frontier-orbitals:hard-soft}{électrophile mou}.

\begin{definition}[Nucléophile ambident]\label{def:b2:frontier-orbitals:ambident}
Un \emph{nucléophile ambident}\index{nucleophile ambident@nucléophile ambident} possède deux atomes nucléophiles
reliés par conjugaison, capables chacun d'attaquer un électrophile.
\end{definition}

L'ion cyanure est attaqué par l'azote, où la charge négative est la plus grande, par les \omterm{def:b2:frontier-orbitals:hard-soft}{électrophiles
durs}, et par le carbone, où sa \omterm{def:b2:frontier-orbitals:frontier}{HOMO} a son plus grand coefficient, par les \omterm{def:b2:frontier-orbitals:hard-soft}{électrophiles mous}~:
un halogénoalcane donne un nitrile \ce{R-C#N}. Les ions énolate, rencontrés au
\cref{ch:b2:enolates-aldol}, réagissent de même par le carbone ou par l'oxygène.

\begin{method}[Une analyse par les orbitales frontières]\label{met:b2:frontier-orbitals:fmo}
\begin{enumerate}
  \item Identifier le nucléophile (\omterm{def:b2:frontier-orbitals:frontier}{HOMO} haute) et l'électrophile (\omterm{def:b2:frontier-orbitals:frontier}{LUMO} basse).
  \item Représenter les deux \omterm{def:b2:frontier-orbitals:frontier}{orbitales frontières} avec leurs coefficients et leurs signes.
  \item La réaction a lieu là où le recouvrement est le plus grand~: atomes de plus grands
        coefficients, lobes de même signe face à face.
  \item La direction d'approche suit la forme de la \omterm{def:b2:frontier-orbitals:frontier}{LUMO} (perpendiculaire au plan d'un
        groupe \ce{C=O}, à l'arrière d'une liaison \ce{C-X}).
  \item Un écart HOMO--LUMO plus petit signifie une réaction plus rapide (toutes choses égales par ailleurs).
\end{enumerate}
\end{method}

La direction d'attaque découle de la \omterm{def:b2:frontier-orbitals:frontier}{LUMO}. La \omterm{def:b2:frontier-orbitals:frontier}{LUMO} d'un groupe carbonyle est son orbitale
$\pi^*$, plus développée sur le carbone et perpendiculaire au plan~: les nucléophiles approchent
le carbone par-dessus ou par-dessous le plan, et non selon l'axe \ce{C=O}. La \omterm{def:b2:frontier-orbitals:frontier}{LUMO} d'un
halogénoalcane est l'orbitale $\sigma^*$ de la liaison \ce{C-X}, dont le gros lobe sur le carbone
pointe à l'opposé de X~: le nucléophile attaque par l'arrière, et le carbone subit une inversion
--- la stéréochimie de la substitution bimoléculaire du volume de première année.

\section{La réaction de Diels--Alder}

\begin{definition}[Réaction concertée]\label{def:b2:frontier-orbitals:concerted}
Une \emph{réaction concertée}\index{reaction concertee@réaction concertée} forme et rompt toutes ses liaisons en
une seule étape, par un état de transition unique, sans intermédiaire.
\end{definition}

\begin{definition}[Réaction de Diels--Alder]\label{def:b2:frontier-orbitals:diels-alder}
La \emph{réaction de Diels--Alder}\index{reaction de Diels--Alder@réaction de Diels--Alder} est l'addition concertée
d'un \emph{diène}\index{diene@diène} conjugué, dans sa conformation s-cis, sur un alcène ou un
alcyne, le \emph{diénophile}\index{dienophile@diénophile}, qui forme un cycle à six chaînons avec deux
nouvelles liaisons $\sigma$ et une nouvelle liaison $\pi$.
\end{definition}

\begin{omfigure}
\begin{tikzpicture}[font=\small, >={Stealth}, scale=1.25]
  % diene s-cis: C1 (0,1.6) C2 (0.6,2.4) C3 (1.6,2.4) C4 (2.2,1.6); dienophile below
  \draw[thick] (0,1.6) -- (0.6,2.4) -- (1.6,2.4) -- (2.2,1.6);
  \draw[thick] (0.2,1.62) -- (0.65,2.22);
  \draw[thick] (2.0,1.62) -- (1.55,2.22);
  \draw[thick] (0.2,0) -- (2.0,0); \draw[thick] (0.35,0.14) -- (1.85,0.14);
  \node[font=\scriptsize, anchor=east] at (-0.05,1.6) {1};
  \node[font=\scriptsize, anchor=south east] at (0.5,2.42) {2};
  \node[font=\scriptsize, anchor=south] at (1.6,2.42) {3};
  \node[font=\scriptsize, anchor=west] at (2.25,1.6) {4};
  % arrows: dienophile pi -> C1..C(dienophile); C1=C2 pi -> C2-C3; C3=C4 pi -> C4-C
  \draw[->, thick, omProp] (1.1,0.22) .. controls (0.9,0.9) and (0.3,0.9) .. (0.12,1.42);
  \draw[->, thick, omProp] (0.45,1.95) .. controls (0.6,2.75) and (1.0,2.8) .. (1.1,2.48);
  \draw[->, thick, omProp] (1.78,1.92) .. controls (2.5,1.6) and (2.4,0.6) .. (2.05,0.2);
  \node at (3.3,1.2) {$\longrightarrow$};
  \begin{scope}[xshift=4.2cm]
    \draw[thick] (0,1.6) -- (0.6,2.4) -- (1.6,2.4) -- (2.2,1.6) -- (2.0,0) -- (0.2,0) -- cycle;
    \draw[thick] (0.72,2.26) -- (1.48,2.26);
  \end{scope}
\end{tikzpicture}

{\small La \omterm{def:b2:frontier-orbitals:diels-alder}{réaction de Diels--Alder} du butadiène et de l'éthène, écrite avec des flèches courbes~:
trois paires d'électrons se déplacent à la fois, en une étape cyclique et concertée~; deux liaisons $\sigma$
se forment aux extrémités du \omterm{def:b2:frontier-orbitals:diels-alder}{diène}, et la liaison $\pi$ se déplace en son centre.}
\end{omfigure}

L'interaction dominante a lieu entre la \omterm{def:b2:frontier-orbitals:frontier}{HOMO} du \omterm{def:b2:frontier-orbitals:diels-alder}{diène}, $\psi_2$, et la \omterm{def:b2:frontier-orbitals:frontier}{LUMO} du
\omterm{def:b2:frontier-orbitals:diels-alder}{diénophile}, $\pi^*$. Leurs lobes aux extrémités du \omterm{def:b2:frontier-orbitals:diels-alder}{diène} et sur les deux carbones du
\omterm{def:b2:frontier-orbitals:diels-alder}{diénophile} ont des signes concordants lorsque le \omterm{def:b2:frontier-orbitals:diels-alder}{diénophile} approche par une face du \omterm{def:b2:frontier-orbitals:diels-alder}{diène}~:
les deux nouvelles liaisons peuvent se former en même temps, sur la même face de chaque partenaire.

\begin{omfigure}
\begin{tikzpicture}[font=\small]
  \foreach \x/\c in {0/0.602, 1/0.372, 2/-0.372, 3/-0.602} {
    \pic at (\x,2) {omp={1.5*\c}{90}};
  }
  \draw[black!45] (-0.2,2) -- (3.2,2);
  \node[anchor=west] at (3.6,2) {HOMO du diène ($\psi_2$)};
  \pic at (0,0) {omp={-0.85}{90}}; \pic at (3,0) {omp={0.85}{90}};
  \draw[black!45] (-0.2,0) -- (3.2,0);
  \node[anchor=west] at (3.6,0) {LUMO du diénophile ($\pi^*$)};
  \draw[omProp, thick, densely dashed] (0,1.45) -- (0,0.6);
  \draw[omProp, thick, densely dashed] (3,1.45) -- (3,0.6);
  \node[font=\scriptsize, omProp, anchor=east] at (-0.15,1.0) {en phase};
  \node[font=\scriptsize, omProp, anchor=west] at (3.15,1.0) {en phase};
\end{tikzpicture}

{\small \omterm{def:b2:frontier-orbitals:frontier}{Orbitales frontières} du butadiène et de l'éthène, lobes proportionnels à leurs
coefficients de Hückel. Face à face, les lobes terminaux de la \omterm{def:b2:frontier-orbitals:frontier}{HOMO} du \omterm{def:b2:frontier-orbitals:diels-alder}{diène} et
les lobes de la \omterm{def:b2:frontier-orbitals:frontier}{LUMO} du \omterm{def:b2:frontier-orbitals:diels-alder}{diénophile} ont les mêmes signes aux deux extrémités~: les deux
liaisons $\sigma$ peuvent se former ensemble, sur la même face de chaque partenaire.}
\end{omfigure}

\begin{proposition}[Stéréospécificité]\label{prop:b2:frontier-orbitals:da-stereospecific}
La \omterm{def:b2:frontier-orbitals:diels-alder}{réaction de Diels--Alder} est stéréospécifique~: des substituants cis sur le \omterm{def:b2:frontier-orbitals:diels-alder}{diénophile} sont
cis dans le produit, des substituants trans restent trans.
\end{proposition}

\begin{proof}
La réaction est concertée et les deux nouvelles liaisons se forment sur la même face du
\omterm{def:b2:frontier-orbitals:diels-alder}{diénophile}~: ses deux carbones ne tournent jamais l'un par rapport à l'autre, et les positions
relatives de leurs substituants sont transportées dans le cycle.
\end{proof}

\begin{proposition}[Régiosélectivité]\label{prop:b2:frontier-orbitals:da-regio}
Avec des partenaires dissymétriques, le produit majoritaire unit l'atome du \omterm{def:b2:frontier-orbitals:diels-alder}{diène} qui a le
plus grand coefficient dans sa \omterm{def:b2:frontier-orbitals:frontier}{HOMO} à l'atome du \omterm{def:b2:frontier-orbitals:diels-alder}{diénophile} qui a le plus grand
coefficient dans sa \omterm{def:b2:frontier-orbitals:frontier}{LUMO}.
\end{proposition}

\begin{proof}[Argument]
Dans l'état de transition, les deux nouvelles liaisons ne sont pas encore également formées~; la
stabilisation $2\beta^2/\Delta\varepsilon$ est plus grande lorsque les plus grands coefficients se recouvrent ($\beta$
croît avec le produit des coefficients des atomes qui se recouvrent)~: cette
orientation est atteinte à plus basse énergie.
\end{proof}

\begin{omfigure}
\begin{tikzpicture}
\begin{axis}[ybar, width=0.47\linewidth, height=4.6cm, ymin=-0.8, ymax=0.8,
  xtick={0,1,2,3,4}, xticklabels={O,C1,C2,C3,C4}, axis x line=bottom, axis y line=left,
  extra y ticks={0}, extra y tick style={grid=major, grid style={black!50}}, enlarge x limits=0.15,
  tick label style={font=\scriptsize}, ylabel={coefficient}, label style={font=\small},
  title={\small HOMO du 1-méthoxybutadiène}, bar width=8pt]
\addplot[fill=omDef!50, draw=omDef] table[x=atom, y=c] {figdata/out/bachelor-2-huckel-c-diene.dat};
\end{axis}
\end{tikzpicture}\hfill
\begin{tikzpicture}
\begin{axis}[ybar, width=0.47\linewidth, height=4.6cm, ymin=-0.8, ymax=0.8,
  xtick={1,2,3,4}, xticklabels={C$\beta$,C$\alpha$,C(O),O}, axis x line=bottom, axis y line=left,
  extra y ticks={0}, extra y tick style={grid=major, grid style={black!50}}, enlarge x limits=0.15,
  tick label style={font=\scriptsize}, ylabel={coefficient}, label style={font=\small},
  title={\small LUMO du propénal}, bar width=8pt]
\addplot[fill=omThm!45, draw=omThm] table[x=atom, y=c] {figdata/out/bachelor-2-huckel-c-dienophile.dat};
\end{axis}
\end{tikzpicture}

{\small Modèle de Hückel avec paramètres d'hétéroatomes ($\alpha_{\ce{O}} = \alpha + 2\beta$,
$\beta_{\ce{CO}} = 0.8\beta$ pour l'oxygène de l'éther~; $\alpha_{\ce{O}} = \alpha + \beta$, $\beta_{\ce{CO}} =
\beta$ pour l'oxygène du carbonyle~; un modèle). La \omterm{def:b2:frontier-orbitals:frontier}{HOMO} du \omterm{def:b2:frontier-orbitals:diels-alder}{diène} est la plus développée sur C4, l'extrémité
éloignée du groupe méthoxy (0.60 contre 0.51)~; la \omterm{def:b2:frontier-orbitals:frontier}{LUMO} du \omterm{def:b2:frontier-orbitals:diels-alder}{diénophile} est la plus développée sur
le carbone $\beta$ (0.66). Ils se lient l'un à l'autre~: les groupes méthoxy et aldéhyde se retrouvent
sur des carbones voisins du cycle.}
\end{omfigure}

Le groupe donneur élève la \omterm{def:b2:frontier-orbitals:frontier}{HOMO} du \omterm{def:b2:frontier-orbitals:diels-alder}{diène} ($m = 0.48$ au lieu de 0.62) et le groupe attracteur
abaisse la \omterm{def:b2:frontier-orbitals:frontier}{LUMO} du \omterm{def:b2:frontier-orbitals:diels-alder}{diénophile} ($m = -0.35$ au lieu de $-1$)~: l'écart se réduit, et de tels
couples réagissent bien plus vite que le butadiène avec l'éthène. Un \omterm{def:b2:frontier-orbitals:diels-alder}{diène} riche en électrons et un
\omterm{def:b2:frontier-orbitals:diels-alder}{diénophile} pauvre en électrons forment la combinaison typique.

\begin{definition}[Adduits endo et exo]\label{def:b2:frontier-orbitals:endo}
Lorsqu'un \omterm{def:b2:frontier-orbitals:diels-alder}{diène} cyclique s'additionne sur un \omterm{def:b2:frontier-orbitals:diels-alder}{diénophile} portant un substituant insaturé, l'\emph{adduit
endo}\index{adduit endo} a ce substituant orienté vers la double liaison nouvellement
formée (sous le \omterm{def:b2:frontier-orbitals:diels-alder}{diène} dans l'état de transition), l'\emph{adduit
exo}\index{adduit exo} à l'opposé. La \emph{règle endo}\index{regle endo@règle endo} énonce
que l'adduit endo se forme plus vite.
\end{definition}

\begin{proposition}[La règle endo]\label{prop:b2:frontier-orbitals:endo}
Sous contrôle cinétique, l'\omterm{def:b2:frontier-orbitals:endo}{adduit endo} est le produit majoritaire, bien que l'\omterm{def:b2:frontier-orbitals:endo}{adduit exo}
soit souvent le plus stable.
\end{proposition}

\admitted

L'explication usuelle est un recouvrement secondaire, dans l'approche endo, entre les
lobes du système $\pi$ du substituant et ceux des atomes centraux de la \omterm{def:b2:frontier-orbitals:frontier}{HOMO}
du \omterm{def:b2:frontier-orbitals:diels-alder}{diène}~; il abaisse l'état de transition endo sans former de liaison. La règle est
empirique et souffre des exceptions~; le volume de troisième année traite ces réactions par la
symétrie de leurs orbitales.

\begin{omfigure}
\begin{tikzpicture}[font=\small]
  \foreach \xs/\mode in {0/endo, 5/exo} {
    \begin{scope}[xshift=\xs cm]
      % cyclopentadiene seen from the side: C1, C2=C3, C4, CH2 bridge on top
      \draw[thick] (0,2) -- (0.6,2.5) -- (1.8,2.5) -- (2.4,2);
      \draw[thick] (0.75,2.38) -- (1.65,2.38);
      \draw[thick] (0,2) -- (1.2,3.2) -- (2.4,2);
      \node[font=\scriptsize, anchor=south] at (1.2,3.2) {\ce{CH2}};
      % dienophile
      \draw[thick] (0.2,0) -- (2.2,0); \draw[thick] (0.35,0.1) -- (2.05,0.1);
      \draw[omProp, densely dashed, thick] (0,1.9) -- (0.2,0.15);
      \draw[omProp, densely dashed, thick] (2.4,1.9) -- (2.2,0.15);
      \node[anchor=north] at (1.2,-1.3) {\mode};
    \end{scope}
  }
  % endo: carbonyls tucked under the diene
  \draw[thick] (0.2,0) -- (0.75,0.95) (2.2,0) -- (1.65,0.95);
  \node[font=\scriptsize, fill=white, inner sep=1pt] at (0.75,1.1) {C=O};
  \node[font=\scriptsize, fill=white, inner sep=1pt] at (1.65,1.1) {C=O};
  \draw[omThm, densely dotted, thick] (0.75,1.25) -- (0.65,2.35);
  \draw[omThm, densely dotted, thick] (1.65,1.25) -- (1.75,2.35);
  % exo: carbonyls pointing away
  \draw[thick] (5.2,0) -- (4.95,-0.85) (7.2,0) -- (7.45,-0.85);
  \node[font=\scriptsize] at (4.95,-1.0) {C=O};
  \node[font=\scriptsize] at (7.45,-1.0) {C=O};
\end{tikzpicture}

{\small Approches endo et exo de l'anhydride maléique sur le cyclopentadiène, schéma.
Les liaisons en formation sont en tirets. Dans l'approche endo, les groupes carbonyle sont sous le
\omterm{def:b2:frontier-orbitals:diels-alder}{diène}, et leurs orbitales $\pi$ recouvrent ses carbones centraux (pointillés, recouvrement
secondaire)~; dans l'approche exo, ils pointent à l'opposé.}
\end{omfigure}

\begin{method}[Dessiner le produit d'une réaction de Diels--Alder]\label{met:b2:frontier-orbitals:diels-alder}
\begin{enumerate}
  \item Placer le \omterm{def:b2:frontier-orbitals:diels-alder}{diène} dans sa conformation s-cis, le \omterm{def:b2:frontier-orbitals:diels-alder}{diénophile} face à ses extrémités.
  \item Tracer les deux nouvelles liaisons $\sigma$ des extrémités du \omterm{def:b2:frontier-orbitals:diels-alder}{diène} aux carbones du \omterm{def:b2:frontier-orbitals:diels-alder}{diénophile},
        et la nouvelle liaison $\pi$ entre C2 et C3 du \omterm{def:b2:frontier-orbitals:diels-alder}{diène}.
  \item Garder les substituants du \omterm{def:b2:frontier-orbitals:diels-alder}{diénophile} du même côté (cis reste cis).
  \item Pour des partenaires dissymétriques, unir les plus grands coefficients (produits «~ortho~» ou «~para~»
        pour les \omterm{def:b2:frontier-orbitals:diels-alder}{diènes} substitués en 1 ou en 2).
  \item Avec un \omterm{def:b2:frontier-orbitals:diels-alder}{diène} cyclique, placer le substituant insaturé en position endo.
\end{enumerate}
\end{method}

\begin{inthelab}[Le craquage du dicyclopentadiène]
Le cyclopentadiène est vendu sous forme de dimère et régénéré juste avant usage~: le dimère est
chauffé jusqu'à son point d'ébullition (environ \qty{170}{\degreeCelsius}) dans un ballon surmonté d'une
colonne de fractionnement, où il se scinde de nouveau en monomère (réaction de rétro-Diels--Alder),
qui distille à \qty{41}{\degreeCelsius} et est recueilli dans un récipient refroidi
dans la glace. Le monomère se dimérise de nouveau en quelques heures à température ambiante et s'utilise
aussitôt. Les deux composés sont inflammables et nocifs~; on travaille sous hotte.
\end{inthelab}

\begin{safety}{\ghs[5mm]{GHS02}\,\ghs[5mm]{GHS06}\,\ghs[5mm]{GHS08}}
Le cyclopentadiène et son dimère sont très inflammables et toxiques par inhalation~; l'anhydride
maléique est corrosif et sensibilisant respiratoire. On les manipule sous hotte,
avec des gants et une protection oculaire, à l'écart des flammes.
% ledger: ghs:cyclopentadiene, ghs:dicyclopentadiene, ghs:maleic-anhydride, bp:cyclopentadiene, bp:dicyclopentadiene
\end{safety}

\section{Exercices}

\begin{exercise}[$\star$]\label{exo:b2:frontier-orbitals:1}
Donner la \omterm{def:b2:frontier-orbitals:frontier}{HOMO} et la \omterm{def:b2:frontier-orbitals:frontier}{LUMO} (énergies en unités de Hückel) de l'éthène, de l'anion allyle
et du butadiène.
\end{exercise}

\begin{exercise}[$\star$]\label{exo:b2:frontier-orbitals:2}
Classer en dur ou mou~: \ce{OH-}, \ce{I-}, \ce{H+}, un thiolate \ce{RS-}, le carbone de
l'iodométhane, \ce{F-}.
\end{exercise}

\begin{exercise}[$\star$]\label{exo:b2:frontier-orbitals:3}
Lesquels de ces \omterm{def:b2:frontier-orbitals:diels-alder}{diènes} peuvent participer à une \omterm{def:b2:frontier-orbitals:diels-alder}{réaction de Diels--Alder}~: buta-1,3-diène,
cyclopentadiène, penta-1,4-diène, (2\textit{Z},4\textit{Z})-hexa-2,4-diène dans sa conformation
s-cis~?
\end{exercise}

\begin{exercise}[$\star$]\label{exo:b2:frontier-orbitals:4}
Dessiner le produit du butadiène avec le propènenitrile \ce{CH2=CH-C#N}.
\end{exercise}

\begin{exercise}[$\star\star$]\label{exo:b2:frontier-orbitals:5}
Deux orbitales à $-10$ et \qty{-2}{eV} interagissent avec $\beta = \qty{-1.0}{eV}$ (données de l'exercice).
Calculer la stabilisation perturbative de deux électrons, puis la stabilisation exacte.
\end{exercise}

\begin{exercise}[$\star\star$]\label{exo:b2:frontier-orbitals:6}
L'ion cyanure réagit avec l'iodométhane pour donner l'éthanenitrile \ce{CH3-C#N}, mais avec un
carbocation dans un solvant fortement ionisant il donne un peu d'isonitrile \ce{R-N#C}. Expliquer.
\end{exercise}

\begin{exercise}[$\star\star$]\label{exo:b2:frontier-orbitals:7}
À l'aide des coefficients du chapitre, prévoir le produit majoritaire de la réaction du 1-méthoxybutadiène
avec le propénal.
\end{exercise}

\begin{exercise}[$\star\star$]\label{exo:b2:frontier-orbitals:8}
Dessiner l'\omterm{def:b2:frontier-orbitals:endo}{adduit endo} du cyclopentadiène avec le propénoate de méthyle.
\end{exercise}

\begin{exercise}[$\star\star$]\label{exo:b2:frontier-orbitals:9}
Pourquoi l'anhydride maléique réagit-il avec le cyclopentadiène à température ambiante, alors que l'éthène
exige un chauffage sous pression~?
\end{exercise}

\begin{exercise}[$\star\star\star$]\label{exo:b2:frontier-orbitals:10}
Le cyclopentadiène réagit avec le maléate de diméthyle (diester cis) et avec le fumarate de diméthyle
(diester trans). Dessiner les deux produits et dire combien de stéréoisomères chacun donne.
\end{exercise}

\begin{exercise}[$\star\star\star$]\label{exo:b2:frontier-orbitals:11}
La \omterm{def:b2:frontier-orbitals:diels-alder}{réaction de Diels--Alder} a un $\Delta_r H^\circ < 0$ et un $\Delta_r S^\circ < 0$. Expliquer ces signes
et pourquoi la réaction inverse devient favorable à haute température.
\end{exercise}

\begin{exercise}[$\star\star\star$]\label{exo:b2:frontier-orbitals:12}
Dans l'hydrazine \ce{H2N-NH2}, chaque azote porte un doublet non liant. À l'aide de l'interaction à
quatre électrons, expliquer pourquoi la molécule évite la conformation dans laquelle les deux doublets
sont parallèles.
\end{exercise}

\section{Problème~: le craquage du cyclopentadiène}

\begin{problem}[{Problème du week-end --- la dimérisation du cyclopentadiène vue comme une
réaction de Diels--Alder, la thermodynamique du craquage, les orbitales frontières du
cyclopentadiène et de l'anhydride maléique, et la stéréochimie de leur adduit}]
\label{pb:b2:frontier-orbitals:1}
Données~: températures d'ébullition, cyclopentadiène \qty{41}{\degreeCelsius}, dicyclopentadiène environ
\qty{170}{\degreeCelsius}. Énergies du modèle de Hückel (un modèle)~: cyclopentadiène, traité comme
un motif butadiène, \omterm{def:b2:frontier-orbitals:frontier}{HOMO} $\alpha + 0.62\beta$ et \omterm{def:b2:frontier-orbitals:frontier}{LUMO} $\alpha - 0.62\beta$~; anhydride maléique, \omterm{def:b2:frontier-orbitals:frontier}{LUMO}
vers $\alpha - 0.35\beta$.
% ledger: bp:cyclopentadiene, bp:dicyclopentadiene

\medskip
\noindent\textbf{Partie I --- Le dimère.}
\begin{enumerate}
  \item Dessiner le cyclopentadiène et montrer que son motif \omterm{def:b2:frontier-orbitals:diels-alder}{diène} est bloqué en s-cis.
  \item Dans la dimérisation, une molécule joue le rôle de \omterm{def:b2:frontier-orbitals:diels-alder}{diène}. Quel est le rôle de l'autre~?
  \item Dessiner le dimère et numéroter les nouvelles liaisons.
  \item Quel adduit, endo ou exo, se forme à température ambiante, et pourquoi~?
  \item La dimérisation est-elle stéréospécifique~? Expliquer à l'aide du mécanisme.
  \item Pourquoi la dimérisation n'a-t-elle pas besoin de catalyseur~?
\end{enumerate}

\medskip
\noindent\textbf{Partie II --- Thermodynamique du craquage.}
\begin{enumerate}[resume]
  \item Donner le signe du $\Delta_r H^\circ$ de la dimérisation, d'après les liaisons formées et
        rompues.
  \item Donner le signe de $\Delta_r S^\circ$.
  \item Écrire $\Delta_r G^\circ(T)$ et montrer qu'elle change de signe à une température $T_i$.
  \item Pourquoi le dimère est-il stable à température ambiante, et le monomère favorisé à chaud~?
  \item Dans le montage de craquage, pourquoi distiller le monomère au fur et à mesure de sa formation
        fait-il avancer la réaction~?
  \item Pourquoi faut-il garder le monomère au froid et l'utiliser rapidement~?
\end{enumerate}

\medskip
\noindent\textbf{Partie III --- Les \omterm{def:b2:frontier-orbitals:frontier}{orbitales frontières}.}
\begin{enumerate}[resume]
  \item Pour le cyclopentadiène réagissant sur lui-même, calculer l'écart HOMO--LUMO en unités de
        $|\beta|$.
  \item Pour le cyclopentadiène avec l'anhydride maléique, calculer l'écart entre la \omterm{def:b2:frontier-orbitals:frontier}{HOMO} du
        \omterm{def:b2:frontier-orbitals:diels-alder}{diène} et la \omterm{def:b2:frontier-orbitals:frontier}{LUMO} de l'anhydride.
  \item Quel couple réagit le plus vite~? Expliquer.
  \item Pourquoi les deux groupes carbonyle de l'anhydride maléique abaissent-ils sa \omterm{def:b2:frontier-orbitals:frontier}{LUMO}~?
  \item Quelle orbitale de l'anhydride maléique assure le recouvrement secondaire de l'approche
        endo~?
\end{enumerate}

\medskip
\noindent\textbf{Partie IV --- L'adduit avec l'anhydride maléique.}
\begin{enumerate}[resume]
  \item Dessiner l'\omterm{def:b2:frontier-orbitals:endo}{adduit endo}.
  \item Repérer ses stéréocentres.
  \item L'adduit est-il chiral~? (Chercher un plan de symétrie.)
  \item Pourquoi les deux hydrogènes de l'ancien \omterm{def:b2:frontier-orbitals:diels-alder}{diénophile} se retrouvent-ils du même côté du
        cycle~?
  \item À quoi ressemblerait l'\omterm{def:b2:frontier-orbitals:endo}{adduit exo}~?
  \item Les \omterm{def:b2:frontier-orbitals:endo}{adduits endo} et exo peuvent-ils s'interconvertir à température ambiante~? Pourquoi~?
  \item Comment l'anhydride réagirait-il avec l'eau, et quel produit se formerait-il~?
  \item Donner le nombre de stéréocentres de l'\omterm{def:b2:frontier-orbitals:endo}{adduit endo} du cyclopentadiène et de
        l'anhydride maléique.
\end{enumerate}
\end{problem}
